% !TeX root = closed-systems-from-comparison-completeness.tex
% ============================================================
\chapter{Standard Results Used}
\label{app:standard-results}
% ============================================================
\begin{chapterorientation}
\orientationitem{Purpose}{This appendix records external mathematical results used as background in the proof chain.}
\orientationitem{Scope}{The statements below are not new results of the closure program.  They identify the standard theorem families invoked in the main text.}
\orientationitem{Use}{The body of the monograph remains the controlling source for all closure-specific hypotheses and conclusions.}
\end{chapterorientation}

This appendix gives a compact reference ledger for standard results used
in the main text.  It is not a substitute for the cited sources and it
does not introduce additional closure-theoretic hypotheses.  Its role is
to make explicit which parts of the argument are supplied by established
background mathematics and which parts are internal to the closure stack.

\section{Boolean representation and compactness}
\label{app:standard-boolean}
The finite-to-global portions of Chapters~\ref{ch:foundation} and
\ref{ch:bottom} use the standard representation of Boolean algebras by
their ultrafilter spaces and the associated compactness principle for
families of closed conditions with the finite intersection property
\cite{Stone1936,Rudin1987}.
In the notation of the main text, this is the background result that
allows a finite-stage admissibility pattern to be tested through the
Stone space
\[
	\UF(\mathsf B)
\]
and its constrained subspace
\[
	\UF(\mathsf B;\mathsf J).
\]
The closure-specific content is not Stone duality itself.  It is the
identification of which Boolean algebra and which exclusion ideal are
generated by intrinsic comparison data.

\section{Category, quotient, and descent language}
\label{app:standard-category}
The use of objects, morphisms, functoriality, quotient maps, and
factorization through endpoint or orbit data is standard categorical
language \cite{MacLane1998}.  In the main text this language is used
only to record descent properties.  For example, a coherent report
descends through a projection precisely when it is constant on the
fibres of that projection.  The closure-specific work is the proof that
the relevant physical projection is forced by intrinsic comparison
semantics rather than chosen externally.

\section{Smooth manifolds, connections, and curvature}
\label{app:standard-smooth}
The realized differential-geometric vocabulary used after the smooth
interface is standard: smooth manifolds, tangent bundles, vector fields,
connections, curvature, covariant derivatives, and Lorentzian metrics
\cite{Lee2013,KobayashiNomizu1963,Nakahara2003,Wald1984}.  The
monograph does not treat these objects as primitive input.  They enter
only after the intrinsic stack has produced the transport data and after
the relevant realization conditions have been stated.

Thus statements such as ``the realized curvature tensor'' or ``the
Levi--Civita connection'' should be read as statements inside the
realized branch.  The internal theorem is the reconstruction of the
transport carrier and its curvature-reading interface; the background
mathematics supplies the standard smooth notation used to read that
interface.

\section{Frame transitivity, constant curvature, and space forms}
\label{app:standard-spaceforms}
The frame-bundle and constant-curvature steps use standard
differential-geometric rigidity results: transitivity of the relevant
frame action forces local isotropy of curvature, and the constant
sectional-curvature models are the familiar space forms
\cite{KobayashiNomizu1963,Wald1984}.  In dimension three, the compact
simply connected positive-curvature space form is the spherical model.
The topological completion of the three-dimensional simply connected
compact case is the Poincare-Perelman theorem, proved through the Ricci
flow program initiated by Hamilton and completed by Perelman
\cite{Hamilton1982,Perelman2002,Perelman2003Ricci,Perelman2003Finite}.

The closure-specific claim is not the space-form classification itself.
The claim is that, at the faithful smooth-realization stage and under
the stated frame-completeness and closure hypotheses, the internal
geometry is driven into the branch to which the standard classification
applies.

\section{Lie and transformation-group structure}
\label{app:standard-lie}
The group-structure trichotomy used in Chapter~\ref{ch:interface} relies
on the standard structure theory of locally compact and transformation
groups: the no-small-subgroups direction of Hilbert's fifth problem,
the Gleason--Yamabe theorem, the structure of totally disconnected
locally compact groups, and the action-side status of the
Hilbert--Smith problem in the regularity regimes cited in the main text
\cite{Gleason1952,Yamabe1953,MontgomeryZippin1955,vanDantzig1936,Tao2014,RepovsScepin1997,Pardon2013}.
These results delimit which ambient group structures can appear once
the intrinsic transport completion is read topologically.

\section{Nilpotent, free, and filtered transport algebra}
\label{app:standard-filtered}
The use of lower central series, graded quotients, and free or partially
commutative Lie-algebra comparisons is standard in group and Lie theory
\cite{Hall1959,Malcev1949,Magnus1935,Witt1937,Labute1970,DuchampKrob1992}.
The main text uses these tools to organize the transport filtration
\[
	F^1\supseteq F^2\supseteq F^3\supseteq\cdots
\]
and to identify which graded piece first carries a stable obstruction.
The internal conclusion is the location and interpretation of the first
closure-visible transport residue.

\section{Hilbert, phase, and operator background}
\label{app:standard-hilbert}
When the text compares the quadratic transport carrier with Hilbert
space language, it uses the standard complex Hilbert-space formalism and
standard perturbation-theoretic conventions
\cite{VonNeumann1955,Kato1995}.  The comparison does not make Hilbert
space primitive.  The monograph's claim is the reconstruction of the
oriented real defect double, its complex structure, and the resulting
Hilbert-side quotient after the relevant transport data have been
constructed.

\section{Finite-order, naturality, and Einstein-boundary results}
\label{app:standard-finite-order}
The interface between intrinsic transport and realized local field laws
uses standard finite-order and naturality results: Peetre's theorem on
local operators, the Epstein--Thurston finite-order theorem for natural
bundles, and Lovelock's classification of divergence-free symmetric
natural tensors of the metric in the relevant dimension
\cite{Peetre1959,EpsteinThurston1979,Lovelock1971}.  The main text uses
these results only after the closure stack has supplied the
second-jet carrier and the realized smooth interface.

\section{Piecewise curvature and transport precedent}
\label{app:standard-piecewise}
The comparison between discrete or combinatorial defect and smooth
curvature is supported by the standard convergence of piecewise-flat
curvature measures in the cited Cheeger--Muller--Schrader framework
\cite{CheegerMullerSchrader1984}.  The dynamics-side precedent in which
closure of a deformation algebra recovers the Einstein Hamiltonian is
the Hojman--Kuchar--Teitelboim result \cite{HojmanKucharTeitelboim1976}.
These are comparison results.  The closure stack supplies its own
transport obstruction and its own interface hypotheses.

\section{Algebraic reconstruction and derivations}
\label{app:standard-algebraic-reconstruction}
The state-side reconstruction extension uses standard facts about
derivations of commutative normed algebras and the spectral
characterization of manifolds \cite{SingerWermer1955,Connes2013}.  These
results are invoked only in the marked extension that studies the
regular readable sector.  They are not premises for the earlier closure,
transport, curvature, phase, charge, or numerical results.

\begin{chapterclosing}
\closingitem{Established}{The external theorem families used by the main text have been isolated from the closure-specific claims.}
\closingitem{Carried forward}{The next appendix gives minimal examples and failure modes for the main constructions.}
\end{chapterclosing}
